\documentclass[12pt,a4paper]{article}

% ── Packages ──────────────────────────────────────────────────────────────────
\usepackage[a4paper, margin=1in]{geometry}
\usepackage{graphicx}
\usepackage{booktabs}
\usepackage{array}
\usepackage{tabularx}
\usepackage[hyphens]{url}
\usepackage{hyperref}
\usepackage{xcolor}
\usepackage{titlesec}
\usepackage{fancyhdr}
\usepackage{enumitem}
\usepackage{tcolorbox}
\usepackage{float}
\usepackage{parskip}
\usepackage[T1]{fontenc}
\usepackage{colortbl}
\usepackage{ragged2e}
\usepackage{subcaption}
\usepackage{caption}

% ── URL line-breaking ──────────────────────────────────────────────────────────
\Urlmuskip=0mu plus 2mu
\def\UrlBreaks{\do\/\do-\do_\do.\do=\do~\do\%}

% ── Hyperref ──────────────────────────────────────────────────────────────────
\hypersetup{
    colorlinks=true,
    linkcolor=blue!70!black,
    urlcolor=blue!70!black,
    citecolor=blue!70!black,
    pdftitle={Project Proposal -- AI Job Market Intelligence},
    pdfauthor={Group 8, CS3012},
    bookmarksnumbered=true,
    bookmarksopen=true,
    breaklinks=true
}

% ── Colors ────────────────────────────────────────────────────────────────────
\definecolor{fastblue}{RGB}{0, 56, 117}
\definecolor{accentblue}{RGB}{55, 138, 221}
\definecolor{lightgray}{RGB}{245, 245, 245}
\definecolor{darkgray}{RGB}{80, 80, 80}
\definecolor{successgreen}{RGB}{26, 158, 110}

% ── Section formatting ────────────────────────────────────────────────────────
\titleformat{\section}{\color{fastblue}\large\bfseries}{\thesection.}{0.5em}{}[\color{accentblue}\titlerule]
\titleformat{\subsection}{\color{fastblue}\normalsize\bfseries}{\thesubsection.}{0.5em}{}
\titleformat{\subsubsection}{\color{darkgray}\normalsize\bfseries}{\thesubsubsection.}{0.5em}{}
\titlespacing{\section}{0pt}{12pt}{6pt}
\titlespacing{\subsection}{0pt}{10pt}{4pt}

% ── Header / Footer ───────────────────────────────────────────────────────────
\pagestyle{fancy}
\fancyhf{}
\fancyhead[L]{\small\color{darkgray} CS3012 -- Data Visualization $\mid$ Spring 2026}
\fancyhead[R]{\small\color{darkgray} Group 8 -- Project Proposal}
\fancyfoot[C]{\small\thepage}
\renewcommand{\headrulewidth}{0.4pt}
\renewcommand{\footrulewidth}{0.4pt}

% ── tcolorbox ─────────────────────────────────────────────────────────────────
\tcbuselibrary{skins,breakable}

\newtcolorbox{infobox}[1][]{
  enhanced, breakable,
  colback=blue!4!white, colframe=accentblue,
  boxrule=0.5pt, arc=4pt, left=8pt, right=8pt, top=6pt, bottom=6pt,
  title={#1}, fonttitle=\bfseries\color{fastblue},
  attach boxed title to top left={yshift=-2mm, xshift=6mm},
  boxed title style={colback=white, colframe=accentblue, boxrule=0.5pt, arc=3pt}
}

\newtcolorbox{storybox}{
  enhanced, colback=fastblue!6!white, colframe=fastblue,
  boxrule=1pt, arc=6pt, left=10pt, right=10pt, top=8pt, bottom=8pt,
}

\newtcolorbox{kpibox}[1][]{
  enhanced, colback=successgreen!8!white, colframe=successgreen,
  boxrule=0.5pt, arc=4pt, left=6pt, right=6pt, top=5pt, bottom=5pt,
  title={#1}, fonttitle=\small\bfseries\color{successgreen!70!black}
}

% ── Misc ──────────────────────────────────────────────────────────────────────
\let\cleardoublepage\clearpage
\newcolumntype{W}{>{\RaggedRight\arraybackslash}X}
\newcolumntype{P}[1]{>{\RaggedRight\arraybackslash}p{#1}}

% ══════════════════════════════════════════════════════════════════════════════
\begin{document}

% ── Cover Page ────────────────────────────────────────────────────────────────
\begin{titlepage}
\thispagestyle{empty}
\centering
\vspace*{0.3cm}
\textcolor{fastblue}{\rule{0.6\linewidth}{2pt}}\\[0.2cm]
{\large\bfseries\color{fastblue} FAST National University of Computer \& Emerging Sciences}\\
{\normalsize\color{darkgray} Islamabad Campus}\\[0.1cm]
\textcolor{fastblue}{\rule{0.6\linewidth}{0.5pt}}\\[0.8cm]
{\small\color{darkgray} Fundamentals of Data Visualization -- CS3012 $\mid$ Section A}\\[0.6cm]
{\Huge\bfseries\color{fastblue} Project Proposal}\\[0.4cm]
{\large\color{accentblue}\rule{0.4\linewidth}{0.5pt}}\\[0.4cm]
{\LARGE\bfseries\color{fastblue}
AI Job Market Intelligence:\\[0.2cm]
Salaries, Roles \& Global Hiring\\[0.2cm]
Trends (2020--2025)}\\[0.6cm]
{\normalsize\color{darkgray}\itshape
Interactive Data Storytelling with Tableau -- A Visualization Project}\\[1.2cm]

\begin{tabular}{>{\bfseries\color{fastblue}}l l l}
\toprule
\# & \textbf{Name} & \textbf{Roll Number} \\
\midrule
1 & Muhammad Nouman Hafeez & 21i-0416 \\
2 & Muhammad Asim         & 21i-0852 \\
3 & Sara Jabeen           & 21i-0623 \\
\bottomrule
\end{tabular}

\vspace{0.8cm}
\begin{tabular}{r P{7cm}}
\textbf{\color{fastblue}Group:}        & 8 \\
\textbf{\color{fastblue}Department:}   & Computer Science \\
\textbf{\color{fastblue}Submitted To:} & Dr.\ Atif Mughees \\
\textbf{\color{fastblue}Semester:}     & Spring 2026 \\
\textbf{\color{fastblue}Date:}         & \today \\
\end{tabular}

\vfill
\textcolor{accentblue}{\rule{0.8\linewidth}{0.5pt}}\\
{\small\color{darkgray} FAST-NUCES Islamabad $\cdot$ CS3012 Data Visualization $\cdot$ Spring 2026}
\end{titlepage}

% ── Table of Contents ─────────────────────────────────────────────────────────
\newpage
\thispagestyle{fancy}
\tableofcontents
\newpage

% ══════════════════════════════════════════════════════════════════════════════
\section{Project Overview}

\subsection{Project Title}
\begin{storybox}
\centering
{\large\bfseries\color{fastblue}
``AI Job Market Intelligence: Salaries, Roles \& Global Hiring Trends (2020--2025)''}
\end{storybox}

\subsection{Short Description}

This project analyzes the global Artificial Intelligence job market using a real-world
dataset of 151{,}445 verified salary submissions sourced from \textbf{aijobs.net} and
its parent platform \textbf{foorilla.com} -- a live salary index covering AI, ML,
Data Science, and Big Data professionals. Through interactive dashboards built in
Tableau Public, the project focuses on AI and ML-relevant roles, uncovering which AI roles command the highest salaries, how
compensation scales with experience level, how remote work adoption has evolved
dramatically from 2020 to 2025, and which countries lead global AI hiring. The
dashboard transforms five years of real anonymous salary data into a compelling,
evidence-backed data story about how AI is reshaping careers, compensation structures,
and the global workforce landscape.

\subsection{Domain of Analytics}

\begin{infobox}[Domain]
\textbf{Primary Domain:} Human Resource \& Workforce Analytics\\[4pt]
\textbf{Sub-Domains:}
\begin{itemize}[leftmargin=1.5em, itemsep=2pt]
  \item \textbf{Compensation Analytics} -- Which AI roles and company sizes offer the highest salaries?
  \item \textbf{Workforce Distribution Analytics} -- How is the AI workforce split across Remote, Hybrid, and On-site?
  \item \textbf{Career Progression Analytics} -- How does salary grow from Entry to Executive?
  \item \textbf{Geographic Analytics} -- Which countries lead global AI hiring and compensation?
  \item \textbf{Temporal Analytics} -- How have salaries and remote work ratios changed year-over-year from 2020 to 2025?
\end{itemize}
\end{infobox}

% ══════════════════════════════════════════════════════════════════════════════
\section{Dataset}

\subsection{Dataset Source \& Platform Evidence}

The salary dataset used in this project originates from
\textbf{aijobs.net} (\url{https://aijobs.net}), a fully live AI/ML job
platform, and is maintained by its parent company \textbf{foorilla.com}
(\url{https://foorilla.com}). The four screenshots below confirm the
platform is real, actively used, and generates its salary data from
genuine live job postings and professional submissions.

\begin{figure}[H]
\centering
\begin{subfigure}[t]{0.48\textwidth}
  \centering
  \includegraphics[width=\linewidth]{pic1.png}
  \caption{foorilla.com Hiring board -- 274{,}213 live job listings
           with real salary ranges, experience levels, and locations.
           This is the live feed that the salary dataset is derived from.}
\end{subfigure}
\hfill
\begin{subfigure}[t]{0.48\textwidth}
  \centering
  \includegraphics[width=\linewidth]{pic2.png}
  \caption{foorilla.com Media feed -- 36{,}866 real-time tech
           articles aggregated from Financial Times, NYT, TechRepublic,
           Reddit and others, confirming an active, multi-source platform.}
\end{subfigure}

\vspace{0.8em}

\begin{subfigure}[t]{0.48\textwidth}
  \centering
  \includegraphics[width=\linewidth]{pic3.png}
  \caption{foorilla.com Insight analytics panel -- 122{,}713 job
           searches, 335{,}339 job views, and 126{,}885 job clicks in the
           past 90 days, proving substantial real user traffic.}
\end{subfigure}
\hfill
\begin{subfigure}[t]{0.48\textwidth}
  \centering
  \includegraphics[width=\linewidth]{pic4.png}
  \caption{aijobs.net job search page -- 32{,}176 live results with
           filters for topic, skill, role, region, experience, and salary.
           The ``Export CSV / JSON'' button directly generates the
           \texttt{salaries.csv} dataset used in this project.}
\end{subfigure}
\caption{Source platform verification -- aijobs.net and foorilla.com are live, active platforms generating real salary data from genuine job postings and professional submissions.}
\end{figure}

\subsection{Dataset Reference}

\begin{table}[H]
\centering
\caption{Dataset Reference Details}
\renewcommand{\arraystretch}{1.4}
\begin{tabularx}{\linewidth}{>{\bfseries\color{fastblue}}l W}
\toprule
Field & Details \\
\midrule
Dataset Name    & Global Salaries in AI, ML, Data Science \& Big Data (aijobs.net Salary Index) \\
Primary Source  & aijobs.net -- anonymous professional salary survey \& live job listings \\
Parent Platform & foorilla.com -- real-time job aggregation and analytics platform \\
Repository      & github.com/foorilla/ai-jobs-net-salaries \\
File Name       & \texttt{salaries.csv} \\
Total Rows      & 151{,}445 verified salary submissions \\
Total Columns   & 11 variables \\
Time Period     & 2020 -- 2025 (five full years) \\
Update Freq.    & Weekly (live dataset, continuously updated) \\
Dataset GitHub  & \url{https://github.com/foorilla/ai-jobs-net-salaries} \\
Kaggle URL      & \url{https://www.kaggle.com/datasets/aijobs/global-salaries-in-ai-ml-data-science} \\
Platform URL    & \url{https://aijobs.net/salaries} \\
Project Repo    & \url{https://github.com/noumanic/ai-job-market-analytics} \\
License         & CC0 1.0 Public Domain (free for any use, no restrictions) \\
\bottomrule
\end{tabularx}
\end{table}

\subsection{Why This Dataset Is Trustworthy}

Unlike many Kaggle datasets that are synthetically generated, the aijobs.net salary
index is collected directly from real AI and ML professionals through an anonymous
self-reporting survey, cross-referenced with live job postings on foorilla.com. As
confirmed by the platform screenshots above, foorilla.com processes over 122{,}000 job
searches and 335{,}000 job views every 90 days, and aijobs.net currently lists
32{,}176 live job openings. The dataset is maintained with weekly updates and published
under the CC0 public domain license with full methodological transparency.

Importantly, the dataset is also published on \textbf{Kaggle} under the official
\texttt{aijobs} account -- confirming it is the authoritative, first-party release
by the platform itself, not a third-party re-upload. The fact that the dataset
publisher and data collector are the same organization is a key credibility indicator.

The dataset exhibits all hallmarks of genuine crowdsourced data: the United States
dominates with over 89\% of submissions -- as expected in a real tech labor market --
full-time employment accounts for 99.4\% of entries, salary distributions are
right-skewed (skewness $= 1.35$) with realistic outliers, and 422 unique job title
variants reflect real-world naming variation. These characteristics stand in sharp
contrast to artificially generated datasets, which typically show perfectly uniform
category distributions and suspiciously few unique values.

\subsection{Key Variables}

\begin{table}[H]
\centering
\caption{Dataset Column Reference}
\renewcommand{\arraystretch}{1.3}
\begin{tabularx}{\linewidth}{>{\ttfamily}l l W}
\toprule
\textbf{Column} & \textbf{Type} & \textbf{Description \& Use in Dashboard} \\
\midrule
work\_year          & Numeric  & Year of submission (2020--2025) -- time-series axis \\
experience\_level   & Category & EN / MI / SE / EX -- career ladder chart and KPI card \\
employment\_type    & Category & FT / PT / CT / FL -- filter dimension (FT = 99.4\%) \\
job\_title          & Text     & Role name (422 unique values) -- grouped for bar chart \\
salary              & Numeric  & Original salary in local currency \\
salary\_currency    & Text     & ISO currency code (USD = 94.5\% of records) \\
salary\_in\_usd     & Numeric  & Standardized salary in USD -- primary measure across all charts \\
employee\_residence & Text     & Country of employee -- geographic filter \\
remote\_ratio       & Numeric  & 0 = On-site, 50 = Hybrid, 100 = Remote -- donut and trend charts \\
company\_location   & Text     & Country of company (97 unique) -- choropleth world map \\
company\_size       & Category & S / M / L -- filter and dual-axis comparison \\
\bottomrule
\end{tabularx}
\end{table}

\subsection{Data Quality Assessment}

\begin{kpibox}[Dataset Health]
\begin{tabularx}{\linewidth}{l l W}
Total Records:     & 151{,}445           & \textcolor{successgreen}{\textbf{Excellent -- largest open AI salary dataset}} \\
Missing Values:    & 0\%                  & \textcolor{successgreen}{\textbf{Perfectly clean, no imputation needed}} \\
Date Range:        & 5 years (2020--2025) & \textcolor{successgreen}{\textbf{Rich longitudinal coverage}} \\
Unique Job Titles: & 422                  & \textcolor{successgreen}{\textbf{Real-world naming diversity}} \\
Unique Countries:  & 97                   & \textcolor{successgreen}{\textbf{Genuine global spread}} \\
Currency:          & Standardized to USD  & \textcolor{successgreen}{\textbf{Normalized via period exchange rates}} \\
Data Origin:       & Live platform survey & \textcolor{successgreen}{\textbf{Verified real professional submissions}} \\
License:           & CC0 Public Domain    & \textcolor{successgreen}{\textbf{No restrictions on academic use}} \\
\end{tabularx}
\end{kpibox}

% ══════════════════════════════════════════════════════════════════════════════
\section{Problem Statement}

The global demand for Artificial Intelligence professionals has surged dramatically since
2020. However, despite widespread discourse about the ``AI job boom,'' there is a lack of
clear, visual, data-driven understanding of \emph{where} AI salaries are highest,
\emph{how much} different roles actually pay across career levels, and \emph{how} work
arrangements have evolved over five years. Students, professionals, and organizations
making career or hiring decisions often rely on anecdotal reports rather than empirical
evidence.

This project addresses that gap by providing a structured, interactive visual analysis
of 151{,}445 real anonymous salary submissions from 2020--2025. Every data point
represents a real AI or ML professional's reported compensation, making the findings
credible, defensible, and directly applicable to career and hiring decisions.

% ══════════════════════════════════════════════════════════════════════════════
\section{Project Objectives}

\begin{enumerate}[leftmargin=2em, itemsep=6pt]
  \item \textbf{Visualize salary distributions} across 15 AI job title categories,
        four experience levels, and three company sizes to identify compensation
        patterns and outliers in the real market.

  \item \textbf{Track temporal salary growth} through a year-over-year trend from
        2020 to 2025, showing how AI compensation has matured over five years.

  \item \textbf{Map the global geography} of AI hiring to reveal which countries lead
        in submissions and average compensation, using real data from 97 countries.

  \item \textbf{Analyze the remote work collapse} year-over-year to quantify the
        dramatic post-pandemic return-to-office shift from 53\% remote in 2022 to
        only 20\% in 2025.

  \item \textbf{Build an interactive Tableau dashboard} with global filters for year,
        experience level, company size, and employment type that apply simultaneously
        across all sheets.

  \item \textbf{Tell a coherent data story} using annotations, KPI cards, and callouts
        that guide viewers through the key findings of the real AI salary landscape.
\end{enumerate}

% ══════════════════════════════════════════════════════════════════════════════
\section{Central Data Story}

\begin{storybox}
\centering
{\normalsize\bfseries\color{fastblue}
``Five Years of Real AI Salaries: A Maturing Market, Widening Pay Gaps,}\\[4pt]
{\normalsize\bfseries\color{fastblue}
and a Dramatic Return to the Office.''}\\[10pt]
\raggedright\normalsize\color{black}
Based on 151{,}445 real salary submissions from 2020--2025, this dashboard reveals that
AI careers command a median compensation of \textbf{\$147{,}000}, that average salaries
grew \textbf{14.7\%} from 2022 to 2024 (\$137{,}747 $\to$ \$157{,}979), that
Executive-level professionals earn nearly \textbf{twice} the salary of Entry-level peers
(\$195{,}896 vs.\ \$102{,}430), and that remote work in AI -- which peaked at 53\% in
2022 -- has sharply declined to just \textbf{20\%} in 2025, signaling a strong
industry-wide return-to-office movement. The United States dominates global AI hiring
with over 89\% of all submissions, and ML Engineers top the compensation ladder at an
average of \textbf{\$195{,}500}.
\end{storybox}

% ══════════════════════════════════════════════════════════════════════════════
\section{Planned Visualizations}

\subsection{Required Charts}

The six required charts are described below, each with its assigned Tableau sheet name,
purpose, and the specific dataset fields used.

\begin{table}[H]
\centering
\caption{Required Chart Plan}
\small % reduces overall size
\renewcommand{\arraystretch}{1.3} % slightly tighter rows

\begin{tabularx}{\linewidth}{
  c
  >{\bfseries\RaggedRight\arraybackslash}p{1.9cm}
  >{\RaggedRight\arraybackslash}p{2.1cm}
  X
}
\toprule
\textbf{\#} & \textbf{Chart Type} & \textbf{Sheet Name} & \textbf{Description \& Fields Used} \\
\midrule

1 & Line Chart
  & Salary Trend
  & Year-over-year avg.\ and median salary (2020--2025).
    \newline\textit{Fields:} \texttt{work\_year}, \texttt{salary\_in\_usd} \\

2 & Bar Chart
  & Salary by Role
  & Roles ranked by avg.\ salary with Top 5/10/All toggle.
    \newline\textit{Fields:} \texttt{job\_title}, \texttt{salary\_in\_usd} \\

3 & Donut Chart
  & Work Mode Split
  & On-site (79\%), Remote (21\%), Hybrid (0.2\%).
    \newline\textit{Fields:} \texttt{remote\_ratio} \\

4 & Map
  & Global AI Hiring
  & World map by avg.\ salary; count in tooltip.
    \newline\textit{Fields:} \texttt{company\_location}, \texttt{salary\_in\_usd} \\

5 & Bar Chart
  & Salary vs Experience
  & Salary across Entry, Mid, Senior, Executive + line overlay.
    \newline\textit{Fields:} \texttt{experience\_level}, \texttt{salary\_in\_usd} \\

6 & KPI Cards
  & Summary KPIs
  & Median, mean, total records, remote \%, EX/EN ratio.
    \newline\textit{Fields:} Aggregated measures \\

\bottomrule
\end{tabularx}
\end{table}

\subsection{Bonus Visualizations}

\begin{table}[H]
\centering
\caption{Bonus Features Plan}
\small
\renewcommand{\arraystretch}{1.3}

\begin{tabularx}{\linewidth}{
  >{\bfseries\RaggedRight\arraybackslash}p{2.4cm}
  >{\RaggedRight\arraybackslash}p{2.3cm}
  X
}
\toprule
\textbf{Bonus Feature} & \textbf{Applied On} & \textbf{Implementation} \\
\midrule

Dual-Axis Chart
  & Salary Trend (Sheet 1)
  & Count as bars (secondary axis) + avg.\ salary line (primary axis);
    shows volume and salary growth together. \\

Global Filters
  & All Sheets
  & Four dropdowns (Year, Experience, Company Size, Employment Type)
    linked across all sheets via Tableau actions. \\

Highlight Actions
  & All Sheets
  & Clicking any element highlights related data across all charts. \\

Parameter Toggle
  & Salary Bar Chart
  & Switch between Top 5, Top 10, and All roles using a parameter
    with calculated ranking. \\

Stacked Bar
  & Remote Trend (Sheet 7)
  & 100\% stacked bars for On-site, Remote, Hybrid showing decline
    from 53\% (2022) to 20\% (2025). \\

\bottomrule
\end{tabularx}
\end{table}

\subsection{Key Annotations Planned}

\begin{itemize}[leftmargin=2em, itemsep=6pt]
  \item \textbf{Salary Trend:}\enspace\textit{``Median salary grew from \$102K (2022) to \$149K (2024) -- a 46\% increase in three years.''}
  \item \textbf{Remote Trend:}\enspace\textit{``Remote work peaked at 53\% in 2022. By 2025, 79\% of AI roles are fully on-site.''}
  \item \textbf{Salary vs Experience:}\enspace\textit{``Executive roles earn 1.91$\times$ entry level -- \$196K vs.\ \$102K.''}
  \item \textbf{Map:}\enspace\textit{``United States leads with 133{,}171 submissions (89\%) and a \$161{,}355 average salary.''}
  \item \textbf{Salary by Role:}\enspace\textit{``ML Engineers top the ladder at \$195{,}500 average -- 91\% above Data Analysts.''}
\end{itemize}

% ══════════════════════════════════════════════════════════════════════════════
\section{Dashboard Design Strategy}

\subsection{Layout Plan}

The master dashboard will follow the structured grid layout shown below.
All panels share a consistent blue color theme and respond to four global
filter dropdowns placed in the title bar.

\vspace{0.5em}
\begin{center}
\renewcommand{\arraystretch}{1.8}
\setlength{\tabcolsep}{4pt}
\begin{tabular}{|p{0.22\linewidth}|p{0.22\linewidth}|p{0.22\linewidth}|p{0.22\linewidth}|}
\hline
\multicolumn{4}{|c|}{\cellcolor{fastblue!18}%
  \small\textbf{Title Bar + Filters: Year \quad Experience \quad Company Size \quad Employment Type}} \\
\hline
\multicolumn{1}{|c|}{\cellcolor{accentblue!12}\small KPI: Median Salary} &
\multicolumn{1}{c|}{\cellcolor{accentblue!12}\small KPI: Mean Salary} &
\multicolumn{1}{c|}{\cellcolor{accentblue!12}\small KPI: EX/EN Ratio} &
\multicolumn{1}{c|}{\cellcolor{accentblue!12}\small KPI: Remote \%} \\[4pt]
\hline
\multicolumn{2}{|p{0.46\linewidth}|}{\cellcolor{lightgray}%
  \centering\small Dual-Axis: Salary Trend 2020--2025} &
\multicolumn{2}{p{0.46\linewidth}|}{\cellcolor{lightgray}%
  \centering\small Bar: Salary by Job Role (param toggle)} \\[14pt]
\hline
\multicolumn{1}{|p{0.22\linewidth}|}{\cellcolor{lightgray}%
  \centering\small Donut:\\ Work Mode} &
\multicolumn{1}{p{0.22\linewidth}|}{\cellcolor{lightgray}%
  \centering\small Bar:\\ Salary vs Exp.} &
\multicolumn{2}{p{0.46\linewidth}|}{\cellcolor{lightgray}%
  \centering\small Choropleth Map: Global Hiring} \\[14pt]
\hline
\multicolumn{4}{|p{0.94\linewidth}|}{\cellcolor{lightgray}%
  \centering\small Stacked Bar: Remote Work Trend 2020--2025 (Bonus Sheet)} \\[8pt]
\hline
\end{tabular}
\end{center}
\vspace{0.5em}

\subsection{Design Principles}

\begin{itemize}[leftmargin=2em, itemsep=4pt]
  \item \textbf{Color scheme:} Consistent blue palette (navy \#003875 primary, sky blue \#378ADD accent) with teal highlights for positive trends.
  \item \textbf{Font:} Tableau Benton Sans -- size 12 body, 14 axis labels, 20 dashboard title.
  \item \textbf{Interactivity:} All filters and highlight actions apply across all sheets simultaneously.
  \item \textbf{Tooltips:} Every mark includes a custom formatted tooltip with the field name, formatted value, and a one-sentence contextual note.
  \item \textbf{Canvas size:} Fixed 1200 $\times$ 900\,px for consistent rendering across screens.
\end{itemize}

% ══════════════════════════════════════════════════════════════════════════════
\section{Data Preparation Plan}

Because this is a real crowdsourced dataset with organic messiness, the following
preprocessing steps will be applied before loading into Tableau.

\begin{enumerate}[leftmargin=2em, itemsep=6pt]
  \item \textbf{Year scope:} Analysis will focus on 2022--2025 where data density is
        sufficient for all chart types (147{,}490 rows, 97\% of all data). Years
        2020--2021 combined contain only 293 rows and will be retained in the trend
        line for historical context only, noted in the chart subtitle.

  \item \textbf{Employment type filter:} Only Full-Time (FT) records will be used for
        salary comparisons (150{,}541 of 151{,}445 rows = 99.4\%). Part-time,
        contract, and freelance roles have structurally different salary ranges and
        would distort comparisons if mixed with FT entries.

  \item \textbf{Outlier handling:} The raw salary range spans \$15{,}000 to
        \$800{,}000. Records outside the 1st--99th percentile (\$36{,}842 to
        \$385{,}000) will be excluded from aggregated charts to prevent extreme values
        from distorting means. The full range remains visible in map tooltips.

  \item \textbf{Job title grouping:} The 422 unique raw title strings will be mapped
        to 15 standardized categories using a calculated field in Tableau (e.g.,
        ``ML Engineer,'' ``Machine Learning Eng.,'' ``ML Eng'' $\to$ ``ML Engineer'').
        The original \texttt{job\_title} column is preserved in all tooltips.

  \item \textbf{Currency:} The \texttt{salary\_in\_usd} column is already standardized
        to USD by the dataset maintainers using period-appropriate exchange rates.
        No additional conversion is required.

  \item \textbf{Remote ratio mapping:} The numeric field (0, 50, 100) will be mapped
        to a string dimension (``On-site,'' ``Hybrid,'' ``Remote'') via a Tableau
        calculated field for readable labels and filter values.
\end{enumerate}

% ══════════════════════════════════════════════════════════════════════════════
\section{Tools and Technologies}

\begin{table}[H]
\centering
\caption{Technology Stack}
\renewcommand{\arraystretch}{1.4}
\begin{tabularx}{\linewidth}{>{\bfseries}l P{2.8cm} W}
\toprule
\textbf{Tool} & \textbf{Version / Platform} & \textbf{Purpose} \\
\midrule
Tableau Public  & Free (2024.x)   & Primary visualization and interactive dashboard tool \\
Microsoft Excel & Office 365      & Initial data inspection and verification before Tableau ingestion \\
Python          & 3.11 (optional) & Supplementary data inspection and title normalization if needed \\
GitHub          & github.com      & Version control, preprocessing scripts, Tableau workbook, and all
                                   project deliverables --
                                   \url{https://github.com/noumanic/ai-job-market-tableau-dashboard} \\
Loom            & loom.com (free) & Screen recording for the 2--5 minute narrated video walkthrough \\
LinkedIn        & linkedin.com    & Article publication and project portfolio promotion \\
\bottomrule
\end{tabularx}
\end{table}

% ══════════════════════════════════════════════════════════════════════════════
\section{Deliverables}

\begin{enumerate}[leftmargin=2em, itemsep=6pt]
  \item \textbf{Tableau Workbook (.twbx)} -- Packaged workbook with the complete
        interactive dashboard and all data embedded, ready for offline viewing.

  \item \textbf{Dashboard PDF / Screenshot} -- High-resolution static export of the
        master dashboard for submission and LinkedIn publishing.

  \item \textbf{Written Summary Report} -- A structured report covering:
        \begin{itemize}[itemsep=2pt]
          \item Dataset source, credibility, and preprocessing steps applied
          \item Key findings from the analysis with supporting statistics
          \item How Tableau features were used to generate insights
          \item Challenges encountered and how they were resolved
        \end{itemize}

  \item \textbf{Video Walkthrough (2--5 minutes)} -- A screen-recorded narrated tour
        of the dashboard demonstrating all charts, filters, annotations, and bonus features.

  \item \textbf{This Project Proposal} -- Formal pre-development documentation of the
        project plan, dataset, objectives, and design strategy.
\end{enumerate}

% ══════════════════════════════════════════════════════════════════════════════
\section{Proposed Timeline}

\begin{table}[H]
\centering
\caption{Project Timeline}
\renewcommand{\arraystretch}{1.5}
\begin{tabularx}{\linewidth}{c >{\bfseries\RaggedRight\arraybackslash}p{3.5cm} W}
\toprule
\textbf{Week} & \textbf{Phase} & \textbf{Tasks} \\
\midrule
1 & Data Preparation
  & Download \texttt{salaries.csv} from GitHub; inspect in Excel; verify 0 missing
    values; identify and document outliers; filter to FT records; confirm currency
    standardization. \\

2 & Tableau Setup \& Exploration
  & Connect CSV to Tableau; define all calculated fields (title grouping, remote
    label, experience label, rank parameter); explore salary distributions by year
    and role. \\

3 & Chart Development (Part 1)
  & Build Sheet 1 (dual-axis salary trend), Sheet 2 (salary by role with parameter
    toggle), Sheet 3 (work mode donut), Sheet 4 (choropleth world map). \\

4 & Chart Development (Part 2)
  & Build Sheet 5 (salary vs experience with overlay), Sheet 6 (KPI cards),
    Sheet 7 (bonus remote work trend stacked bar). \\

5 & Dashboard Assembly
  & Assemble master dashboard; configure all global filter actions and highlight
    actions; add title bar, KPI strip, color theme, and all five text annotations. \\

6 & Storytelling \& Polish
  & Refine all custom tooltips; finalize parameter toggle behavior; adjust fonts,
    spacing, and color consistency across all sheets. \\

7 & Submission \& Presentation
  & Export \texttt{.twbx}; record Loom walkthrough video; write summary report;
    compile and submit final package to Dr.\ Atif Mughees. \\
\bottomrule
\end{tabularx}
\end{table}

% ══════════════════════════════════════════════════════════════════════════════
\section{Future Significance}

This project holds significance beyond the classroom in several dimensions:

\begin{enumerate}[leftmargin=2em, itemsep=6pt]

  \item \textbf{Career Guidance for CS Students:} The dashboard provides Computer
        Science graduates with evidence-based insight into which AI specializations
        command the highest compensation, enabling more informed decisions about
        which roles and skills to prioritize.

  \item \textbf{Return-to-Office Trend Awareness:} The sharp decline in remote work
        from 53\% in 2022 to only 20\% in 2025 is a critical signal for anyone
        planning a career around location flexibility. This dashboard makes that
        five-year trend visible and filterable.

  \item \textbf{Salary Negotiation Benchmark:} With 151{,}445 real data points, the
        dashboard serves as a credible benchmarking tool for AI/ML professionals in
        salary negotiations, with breakdowns by role, experience, and country.

  \item \textbf{Hiring Strategy for Organizations:} HR departments can use the
        geographic and role-level insights to set competitive compensation benchmarks
        and identify where AI talent is most concentrated globally.

  \item \textbf{Portfolio Demonstration:} This project demonstrates real-world data
        literacy, Tableau proficiency, and end-to-end storytelling ability using a
        live, continuously updated industry dataset -- skills directly valued by data
        analyst, BI, and data science employers.

  \item \textbf{Scalability:} Because the aijobs.net / foorilla.com dataset is updated
        weekly, the dashboard can be refreshed with new data at any time, enabling
        ongoing year-over-year trend analysis as the AI job market continues to evolve.

\end{enumerate}

% ══════════════════════════════════════════════════════════════════════════════
\section{References}

\begin{enumerate}[leftmargin=2em, itemsep=8pt]

  \item aijobs.net. (2025). \textit{Global Salaries in AI, ML and Data Science}
        [Data set]. CC0 1.0 Public Domain. foorilla/ai-jobs-net-salaries.
        \\\url{https://github.com/foorilla/ai-jobs-net-salaries}

  \item aijobs.net. (2025). \textit{Global Salaries in AI, ML, Data Science}
        [Official Kaggle dataset, published under \texttt{aijobs} account].
        \\\url{https://www.kaggle.com/datasets/aijobs/global-salaries-in-ai-ml-data-science}

  \item aijobs.net. (2025). \textit{AI, ML and Data Science Job Search Platform}.
        \\\url{https://aijobs.net}

  \item foorilla.com. (2025). \textit{foorilla -- Real-Time Job Aggregation and
        Talent Analytics Platform}.
        \\\url{https://foorilla.com}

  \item Tableau Software. (2024). \textit{Tableau Public -- Free Data Visualization Tool}.
        \\\url{https://public.tableau.com}

  \item Tableau Software. (2024). \textit{Dashboard Design Best Practices}.
        Tableau Help Documentation.
        \\\url{https://help.tableau.com/current/pro/desktop/en-us/dashboards_best_practices.htm}

\end{enumerate}

\vspace{1cm}
\begin{center}
\textcolor{accentblue}{\rule{0.6\linewidth}{0.5pt}}\\[4pt]
{\small\color{darkgray}
Group 8 $\cdot$ CS3012 Fundamentals of Data Visualization $\cdot$ Spring 2026\\
FAST-NUCES Islamabad $\cdot$ Submitted to Dr.\ Atif Mughees}
\end{center}

\end{document}